\documentclass[11pt,a4paper]{report}

% ==========================================
% NEW ELEGANT MINIMAL B/W PREAMBLE
% ==========================================
\usepackage[utf8]{inputenc}
\usepackage[T1]{fontenc}
\usepackage[german]{babel}
\usepackage{amsmath, amssymb, amsthm, mathtools, environ}
\usepackage{booktabs} % Required for \toprule, \midrule, and \bottomrule
\usepackage{fontawesome5} % For the environment symbols

% --- FONT CONFIGURATION ---
\usepackage{mathpazo} % Elegant Palatino font

% --- EMPHASIS STYLE OVERRIDE ---
\DeclareTextFontCommand{\emph}{\textbf} % Redefine \emph to use \textbf

\usepackage[margin=1in, headheight=35pt]{geometry}
\usepackage[dvipsnames]{xcolor} % Keeping xcolor for grayscale definitions
\usepackage{hyperref}
\usepackage[nameinlink, capitalize]{cleveref}

% --- HYPERREF AESTHETICS ---
\hypersetup{
    colorlinks=true,
    linkcolor=black,
    filecolor=darkgray,      
    urlcolor=black,
    citecolor=black,
    bookmarksopen=true
}

\usepackage{tcolorbox}
\tcbuselibrary{skins, breakable, theorems}

% --- TYPOGRAPHY & LAYOUT ENHANCEMENTS ---
\usepackage{parskip} 
\linespread{1.05}    

% --- LIST SPACING FIX ---
\usepackage{enumitem}
\setlist[enumerate]{
    topsep=0.5ex,
    itemsep=0.5ex,
    parsep=\parskip,
    partopsep=0pt,
    label=\textbf{(\alph*)},
    ref=(\alph*)
}

\setlist[description]{
    style=unboxed,
    labelsep=0.333em plus 0.167em minus 0.111em
}

\usepackage{fancyhdr}
\usepackage{extramarks}
\pagestyle{fancy}
\fancyhf{}
\fancyhead[L]{\sffamily\color{black}\makebox[0.4\textwidth][l]{\leftmark}}
\fancyhead[R]{\sffamily\color{black}\makebox[0.4\textwidth][r]{\thepage}}
\fancyfoot[L]{\sffamily\color{black}\makebox[0.4\textwidth][l]{\rightmark}}
\fancyfoot[R]{\sffamily\color{black}\makebox[0.4\textwidth][r]{\lastxmark}}
\renewcommand{\headrule}{\hbox to\headwidth{\color{gray!50}\leaders\hrule height 0.5pt\hfill}}
\renewcommand{\footrule}{\hbox to\headwidth{\color{gray!50}\leaders\hrule height 0.5pt\hfill}}

\renewcommand{\sectionmark}[1]{%
    \markright{\MakeUppercase{\thesection.\ #1}}%
    \extramarks{}{}%
}
\renewcommand{\subsectionmark}[1]{%
    \extramarks{\MakeUppercase{\thesubsection.\ #1}}{\MakeUppercase{\thesubsection.\ #1}}%
}

% Fail-Safe TikZ Configuration
\usepackage{tikz}
\usetikzlibrary{calc, arrows.meta, positioning, 3d, shapes.geometric, decorations.pathreplacing, patterns}
\usepackage{pgfplots}
\pgfplotsset{compat=1.18}

\usepackage{titlesec}

\definecolor{bglight}{HTML}{F8F9FA}    

% Math operators
\DeclareMathOperator{\id}{id}
\DeclareMathOperator{\Vol}{Vol}
\newcommand{\newterm}[1]{\emph{#1}}
\newcommand{\inlinemetanote}[1]{\textit{[\textcolor{gray}{#1}]}}

\renewcommand{\textbf}[1]{{\bfseries\sffamily#1}}

% Custom Theorem Style
\newtheoremstyle{bwthm}%
  {\topsep}{\topsep}{\normalfont}{}{\sffamily\bfseries\color{black}}{:}{0.5em}%
  {\thmname{#1}\thmnumber{ #2}\thmnote{ \textbf{(}{\normalfont\sffamily\textit{#3}}\textbf{)}}}

\theoremstyle{bwthm}

% Standard Math Environments
\newtheorem{theorem}{Theorem}[chapter]
\newtheorem{lemma}[theorem]{Lemma}
\newtheorem{definition}[theorem]{Definition}
\newtheorem{corollary}[theorem]{Corollary}
\newtheorem{proposition}[theorem]{Proposition}
\newtheorem{example}[theorem]{Example}
\newtheorem{exercise}[theorem]{Exercise}

% Unnumbered Environments
\newtheorem*{theorem*}{Theorem}
\newtheorem*{lemma*}{Lemma}
\newtheorem*{proposition*}{Proposition}
\newtheorem*{definition*}{Definition}
\newtheorem*{corollary*}{Corollary}
\newtheorem*{claim*}{Claim}
\newtheorem*{example*}{Example}
\newtheorem*{exercise*}{Exercise}

\newtheorem*{notation}{Notation}
\newtheorem*{remark}{Remark}
\newtheorem*{summary}{Summary}
\newtheorem*{warmup}{Warm up}
\newtheorem*{question}{Question}
\newcommand{\qt}[1]{\textit{``#1''}}
\newtheorem*{answer}{Answer}
\newtheorem*{importantremark}{Important remark}
\newtheorem*{goals}{Goals}

\crefalias{lemma}{lemma}
\crefalias{theorem}{theorem}
\crefalias{definition}{definition}
\crefalias{proposition}{proposition}

% ==========================================
% REDESIGNED MINIMAL B/W ENVIRONMENTS
% ==========================================

% 1. Clean Spoken Protocol
\newtcolorbox{spoken-clean}[1][]{
    enhanced,
    frame hidden,
    colback=gray!5,
    borderline west={3pt}{0pt}{black},
    title={\sffamily\bfseries\textcolor{black}{\faMicrophone\ \ Spoken Protocol [#1]}},
    attach title to upper=\par\vspace{1ex},
    breakable
}

% 2. AI Stroke-by-Stroke Note
\newtcolorbox{nice-box}[1][]{
    enhanced,
    colback=white,
    colframe=black,
    boxrule=0.5pt,
    drop fuzzy shadow,
    title={\sffamily\bfseries\faStickyNote\ \ #1},
    coltitle=black,
    attach boxed title to top left={xshift=10pt, yshift=-10pt},
    boxed title style={colback=gray!20, colframe=black, boxrule=0.5pt},
    top=14pt,
    breakable
}

% 3. Best Teacher / Didactic Insight
\newtcolorbox{didactic-insight}[1][Pedagogical Concept]{
    enhanced,
    arc=3mm,
    frame hidden,
    colback=gray!5,
    borderline west={2pt}{0pt}{gray!80},
    title={\sffamily\bfseries\textcolor{black}{\textit{\faLightbulb\ \ Insight:} #1}},
    attach title to upper=\par\vspace{0.5ex},
    breakable
}

% 4a. Modular Visual Blackboard Replication Box
\NewTColorBox{color-box}{ m O{} }{
    enhanced,
    colframe=black,
    colback=white,
    colupper=black,
    fonttitle=\bfseries\sffamily,
    coltitle=black,
    title={\faBox\ \ #2},
    boxrule=1.5pt,
    arc=4pt,
    drop fuzzy shadow=gray!30,
    attach boxed title to top left={xshift=10pt, yshift=-10pt},
    boxed title style={colback=white, colframe=black, boxrule=0.5pt, top=2pt, bottom=2pt, left=4pt, right=4pt},
    if empty title={notitle}
}

% 4b. Legacy Visual Blackboard Replication Boxes (Backward Compatibility)
\newcommand{\NewChalkBox}[3]{%
    \newtcolorbox{#1}[1][Highlighted Board Formula]{%
        enhanced,
        colback=white,
        colframe=black,
        fonttitle=\bfseries\sffamily,
        coltitle=black,
        title=\faBox\ \ ##1,
        boxrule=1.5pt,
        arc=4pt,
        drop fuzzy shadow=gray!30,
        attach boxed title to top left={xshift=10pt, yshift=-10pt},
        boxed title style={colback=white, colframe=black, boxrule=0.5pt, top=2pt, bottom=2pt, left=4pt, right=4pt}
    }%
}

% Alias all colors to the minimal black box to maintain compilation compatibility
\NewChalkBox{orange-box}{black}{gray!30}
\NewChalkBox{blue-box}{black}{gray!30}
\NewChalkBox{dark-green-box}{black}{gray!30}
\NewChalkBox{apple-green-box}{black}{gray!30}
\NewChalkBox{yellow-box}{black}{gray!30}
\NewChalkBox{violet-box}{black}{gray!30}
\NewChalkBox{red-box}{black}{gray!30}
\NewChalkBox{green-box}{black}{gray!30}
\NewChalkBox{purple-box}{black}{gray!30}
\NewChalkBox{gray-box}{black}{gray!30}
\NewChalkBox{apricot-box}{black}{gray!30}
\NewChalkBox{bittersweet-box}{black}{gray!30}
\NewChalkBox{teal-blue-box}{black}{gray!30}
\NewChalkBox{cyan-box}{black}{gray!30}
\NewChalkBox{magenta-box}{black}{gray!30}
\NewChalkBox{brown-box}{black}{gray!30}
\NewChalkBox{pink-box}{black}{gray!30}

\NewChalkBox{orange-formula}{black}{gray!30}
\NewChalkBox{blue-formula}{black}{gray!30}
\NewChalkBox{dark-green-formula}{black}{gray!30}
\NewChalkBox{apple-green-formula}{black}{gray!30}
\NewChalkBox{yellow-formula}{black}{gray!30}
\NewChalkBox{violet-formula}{black}{gray!30}
\NewChalkBox{red-formula}{black}{gray!30}

% 5. Mathematical Setup (Stroke-by-Stroke)
\newtcolorbox{math-stroke}[1][Mathematical Setup]{
    enhanced,
    frame hidden,
    colback=white,
    opacityback=0,
    borderline west={2pt}{0pt}{gray!40},
    title={\sffamily\bfseries\textcolor{black}{\faChalkboard\ \ #1}},
    attach title to upper=\par\vspace{1ex},
    breakable
}

% 6. Redundant Explanation
\newtcolorbox{redundant-explanation}[1][Redundant Explanation of Step]{
    enhanced,
    arc=2mm,
    colback=gray!15,
    colframe=black,
    title={\sffamily\bfseries\faHistory\ \ #1},
    breakable
}

% 7. Meta-Note
\newtcolorbox{meta-note}[1][Administrative Meta-Note]{
    enhanced,
    frame hidden,
    colback=white,
    borderline={0.5pt}{0pt}{gray, dashed},
    title={\sffamily\bfseries\textcolor{gray}{\faEye\ \ #1}},
    attach title to upper=\par\vspace{1ex},
    breakable
}

% 8. AI Observational Note
\newtcolorbox{ai-note}[1][AI Note]{
    enhanced,
    colback=gray!20,
    colframe=black,
    fonttitle=\bfseries\sffamily,
    coltitle=black,
    title=\faRobot\ \ #1,
    boxrule=1pt,
    arc=3pt,
    breakable
}

% 9. Student Question / Interaction
\newtcolorbox{student-interaction}[1][Student Interaction]{
    enhanced,
    frame hidden,
    colback=gray!10,
    borderline west={3pt}{0pt}{gray!60},
    title={\sffamily\bfseries\textcolor{black}{\faUserGraduate\ \ #1}},
    attach title to upper=\par\vspace{1ex},
    left skip=20pt,
    right skip=20pt,
    breakable
}

% 10. Explanation of Steps
\newenvironment{explanation-of-steps}[1][Explanation of Steps]{
    \par\vspace{1ex}\noindent\textbf{\faListUl\ \ #1:}
}{
    \par\vspace{1ex}
}

% 10.5 Short Inline Proof
\newenvironment{short-proof}[1][Proof]{
    \par\vspace{1ex}\noindent
    \tikz[baseline=(A.base)]\node[fill=gray!80, text=white, rounded corners=2pt, inner xsep=4pt, inner ysep=2pt, font=\sffamily\bfseries\small] (A) {\faPenNib\ \ #1};
    \par\vspace{1ex}\noindent
}{
    \hfill\tikz[baseline=(Q.base)]\node[draw=gray!80, fill=white, line width=0.8pt, rounded corners=1.5pt, inner xsep=3pt, inner ysep=1.5pt, font=\sffamily\bfseries\scriptsize, text=gray!80] (Q) {Q.E.D.};\par\vspace{1ex}
}

% 10.75 Subsection-like heading for proofs
\newenvironment{proof-subsection}[1]{
    \par\vspace{1ex}\noindent
    \tikz[baseline=(A.base)]\node[fill=gray!30, text=black, rounded corners=2pt, inner xsep=4pt, inner ysep=2pt, font=\sffamily\bfseries] (A) {#1};
    \vspace{0.5ex}\par 
}{
    \par\vspace{1ex}
}

% 11. Proof Wrapper
\let\proof\relax
\let\endproof\relax
\newtcolorbox{proof}[1][Proof]{
    enhanced,
    breakable,
    frame hidden,
    colback=white,
    opacityback=0,
    borderline west={3pt}{0pt}{gray!80},
    title={\tikz[baseline=(A.base)]\node[fill=gray!80, text=white, rounded corners=2pt, inner xsep=6pt, inner ysep=3pt, font=\sffamily\bfseries\large] (A) {\faPenNib\ \ #1};},
    attach title to upper=\par\vspace{1.5ex},
    finish unbroken={%
        \node[rotate=35, scale=5, text=gray, opacity=0.12, font=\sffamily\bfseries] at (frame.center) {PROOF};
        \node[anchor=south east, xshift=-4pt, yshift=4pt] at (frame.south east) {\tikz\node[draw=gray!80, fill=white, line width=1pt, rounded corners=2pt, inner xsep=4pt, inner ysep=2pt, font=\sffamily\bfseries\small, text=gray!80] {Q.E.D.};};
    },
    finish first={%
        \node[rotate=35, scale=5, text=gray, opacity=0.12, font=\sffamily\bfseries] at (frame.center) {PROOF};
    },
    finish middle={%
        \node[rotate=35, scale=5, text=gray, opacity=0.12, font=\sffamily\bfseries] at (frame.center) {PROOF};
    },
    finish last={%
        \node[rotate=35, scale=5, text=gray, opacity=0.12, font=\sffamily\bfseries] at (frame.center) {PROOF};
        \node[anchor=south east, xshift=-4pt, yshift=4pt] at (frame.south east) {\tikz\node[draw=gray!80, fill=white, line width=1pt, rounded corners=2pt, inner xsep=4pt, inner ysep=2pt, font=\sffamily\bfseries\small, text=gray!80] {Q.E.D.};};
    }
}

% 12. Lecture Break
\newtcolorbox{lecture-break}[1][Lecture Break]{
    enhanced,
    frame hidden,
    colback=white,
    borderline north={0.5pt}{0pt}{gray, dashed},
    borderline south={0.5pt}{0pt}{gray, dashed},
    title={\sffamily\bfseries\centering\textcolor{gray}{\faCoffee\ \ #1}},
    attach title to upper=\par\vspace{1ex},
    fontupper=\centering\itshape\color{gray},
    breakable
}

% ==========================================
% INVISIBLE LAYER & FALLBACK ENVIRONMENTS
% ==========================================
\NewEnviron{ai-tikz-planner-invisible-content}{}
\NewEnviron{ai-type-check-invisible-content}{}
\NewEnviron{ai-proof-skeleton-invisible-content}{}
\NewEnviron{ai-example-invisible-content}[1][]{}
\NewEnviron{ai-discard-invisible-content}[1][]{}
\NewEnviron{ai-retroactive-patch}[1][]{}
\NewEnviron{ai-global-state-checkpoint-invisible-content}{}

% Data Integrity Fallbacks (Grayscale error boxes)
\newcommand{\NewFallbackBox}[1]{
    \newtcolorbox{#1}[1][Fallback]{
        enhanced,
        colback=gray!15,
        colframe=black,
        title={\sffamily\bfseries\textcolor{white}{\faExclamationTriangle\ \ ##1}},
        fonttitle=\sffamily\bfseries,
        coltitle=black,
        title={##1},
        breakable
    }
}
\NewFallbackBox{ai-logical-gap}
\NewFallbackBox{ai-generation-loop-fallback}
\NewFallbackBox{ai-raw-ocr-fallback}
\NewFallbackBox{ai-phonetic-dump}
\NewFallbackBox{ai-modality-conflict}
\NewFallbackBox{ai-off-camera-state}
\NewFallbackBox{ai-async-board-update}
\NewFallbackBox{ai-kinetic-emphasis}

% Custom Lecture Chapter Command
\newcommand{\lecturechapter}[4]{%
  \chapter[#1, #2: #4]{{\Large\color{black}#1: #3} \texorpdfstring{\\[1ex]}{-- } {\LARGE #4}}%
}

\titlespacing*{\chapter}{0pt}{-20pt}{40pt}

\titleformat{\chapter}[display]
  {\normalfont\sffamily\Large\bfseries\color{black}}
  {\raggedleft\large\color{gray}\MakeUppercase{\chaptertitlename}\ \thechapter}
  {1ex}
  {\color{black}\rule{\textwidth}{1.5pt}\vspace{2ex}\raggedleft \break}
  [\vspace{1ex}\color{black}\rule{\textwidth}{1.5pt}]

\titleformat{\section}
  {\normalfont\sffamily\Large\bfseries\color{black}} 
  {\thesection}{1em}{}

\titleformat{\subsection}
  {\normalfont\sffamily\large\bfseries\color{black}}
  {\thesubsection}{1em}{}

\titleformat{\subsubsection}
  {\normalfont\sffamily\normalsize\bfseries\color{darkgray}}
  {\thesubsubsection}{1em}{}

% Level 1: Section
\newcommand{\sectionbar}{%
  \tikz[baseline=0pt]{
     \fill[black] (0,0) rectangle (3.5pt, 0.72em);
    \fill[darkgray] (4.8pt,0) rectangle (6.3pt, 0.72em);
    \fill[gray] (7.6pt,0) rectangle (8.1pt, 0.72em);
  }%
  \hspace{10pt}%
}

% Level 2: Subsection
\newcommand{\subsectionbar}{%
  \tikz[baseline=0pt]{
    \shade[left color=black, right color=gray] (0,0) rectangle (2.8pt, 0.72em);
    \fill[lightgray] (4.1pt,0) rectangle (4.8pt, 0.72em);
  }%
  \hspace{10pt}%
}

% Level 3: Subsubsection
\newcommand{\subsubsectionbar}{%
  \tikz[baseline=0pt]{
    \fill[gray] (0,0) rectangle (2.2pt, 0.72em);
  }%
  \hspace{10pt}%
}

\titleformat{\section}{\normalfont\sffamily\Large\bfseries\color{black}}{}{0pt}{\makebox[0pt][r]{\sectionbar}}
\titleformat{\subsection}{\normalfont\sffamily\large\bfseries\color{darkgray}}{}{0pt}{\makebox[0pt][r]{\subsectionbar}}
\titleformat{\subsubsection}{\normalfont\sffamily\normalsize\bfseries\color{gray}}{}{0pt}{\makebox[0pt][r]{\subsubsectionbar}}

\titlespacing*{\section}{0pt}{3.5ex plus 1ex minus .2ex}{1.5ex plus .2ex}
\titlespacing*{\subsection}{0pt}{3ex plus 1ex minus .2ex}{1.2ex plus .2ex}
\titlespacing*{\subsubsection}{0pt}{2.5ex plus 1ex minus .2ex}{1ex plus .2ex}
